\documentclass[11pt,a4paper]{article}

\usepackage[utf8]{inputenc}
\usepackage[spanish, es-tabla]{babel}
\usepackage[margin=2.5cm]{geometry}

% Paquetes matemáticos
\usepackage{amsmath}
\usepackage{amsfonts}
\usepackage{amssymb}
\usepackage{amsthm}

% Paquete de código (Indispensable para que compile)
\usepackage{listings}
\usepackage{xcolor}

\lstset{
    backgroundcolor=\color{gray!5},
    basicstyle=\ttfamily\small,
    breaklines=true,
    frame=single,
    tabsize=2
}

\title{Prueba de Entorno Exactas}
\author{Walter Fregenal}
\date{\today}

\begin{document}

\maketitle

\section{Introducción}
Si esto compila bien, el entorno local de \LaTeX{} en Ubuntu está perfectamente integrado con VS Code.

\section{Demostración Matemática}
Una ecuación integrada en el texto: $e^{i\pi} + 1 = 0$.

Y una ecuación destacada:
\begin{equation}
    f(x) = \int_{-\infty}^{\infty} \hat{f}(\xi)\,e^{2\pi i \xi x}\,d\xi
\end{equation}

\section{Prueba de Snippets}
A partir de la línea de abajo podés probar tus atajos mágicos:
\begin{lstlisting}[language=Haskell, caption=${:Descripción}]
main = putStrLn "Hello, World!"
\end{lstlisting}
\end{document}

